\section{Conclusion}
\label{sec:conclusion}

We asked whether a pass on a tagged assessment item requires the tagged competency, and answered it for one subject and one set of public items with executable counterexamples rather than with student data or expert judgment. On New York Regents Algebra~I and Algebra~II the published key is recoverable by a solver that never sees the given quantities, so a score there is read as evidence of an operation on inputs the solver did not have. On a few other cells it is recoverable by programs denied the tagged operation. On the TIMSS cognitive-domain cells it is recoverable by none of the solvers we tried at the sizes available. Each of those statements is a property of an item and its scoring rule, and each was checked before any student sat the test.

What we built is one instance of an automated curriculum falsifier. We do not expect this instance to be deployed. What it shows is that the inputs the search needs, the tag text, the item text, the keys, an executable scoring rule, and a restriction on the solver that a machine can enforce, are already public for large assessment programs, and that where they are public the search runs. That makes a bypass search a natural member of an assessment test suite. A discovered bypass becomes a regression test that a revised item must fail and a competent reference solver must still pass, run by item writers before release, by purchasers before adoption, and by researchers before a tagged item serves as an outcome measure. The same holds for any subject whose items and scoring rules can be executed, and it reaches constructed response as soon as a rubric has an executable scorer.

The measurement lesson travels with it: the text layer of a released PDF is not the test, and an automated audit of released items has to read the page.
